%% Fiche de préparation de séance (format type INSPE) -- document enseignant.
%% Séance 1 de la séquence 4 (séance observée lors de la visite d'inspection du 13/10/2026).
%% Calcul mental : diaporama Act-6C04-03-Calcul-Mental-Complements.pdf.
%% Activité : fiche Act-6C04-01-Decouverte-des-Angles.pdf. Pas d'ardoise : tout dans le cahier d'exercices.

%% Font size %%
\documentclass[11pt]{article}

%% Load the custom package
\usepackage{Mathdoc}

%% Numéro de séquence %% Titre de la séquence %%
\renewcommand{\centerhead}{S04 - Préparation de séance}

%% Spacing commands %%
\renewcommand{\baselinestretch}{1} \setlength{\parindent}{0pt}

\newcommand{\champ}[2]{\textbf{#1} & #2 \\ \hline}
\newcommand{\minutes}[1]{\textbf{#1 min}}


\begin{document}

\begin{center}
{\Large\bfseries Fiche de préparation de séance}\\[1mm]
{\large Séquence 4 : Les angles -- Séance 1 : découvrir les angles, les nommer}
\end{center}

\renewcommand{\arraystretch}{1.35}
\begin{tabularx}{\linewidth}{|>{\raggedright\arraybackslash}p{3.6cm}|X|}
\hline
\champ{Classe / effectif}{6\textsuperscript{e} D -- 22 élèves (REP+)}
\champ{Date / horaire}{Mardi 13 octobre 2026, 10 h 10 -- 11 h 05 \qquad Durée : 55 min}
\champ{Place dans la séquence}{Séance 1 sur 10 (plan détaillé : \textbf{Seq-6C04-Plan-de-Sequence}).\enskip
\textbf{S1 : découverte (Act-6C04-01), définition, sommet, côtés, nommer un angle.}\enskip
S2 : comparer des angles (gabarits), angles particuliers.\enskip
S3 : évaluation 1.\enskip S4 : mesurer.\enskip S5 : construire.\enskip S6 : évaluation 2.\enskip S7 : salle info, bilan sur Capytale.\enskip S8 : salle info, permis rapporteur.\enskip S9 : évaluation 3.\enskip S10 : QCM bilan.\newline
Séquence précédente en géométrie : S02 (points, droites, demi-droites, perpendiculaires).}
\champ{Domaine du programme}{Cycle 3 -- Espace et géométrie : reconnaître, nommer, décrire des figures ; \textbf{utiliser le vocabulaire et les notations} géométriques (demi-droite, angle, sommet) ; l'angle comme \textbf{ouverture} entre deux demi-droites, première étape vers la comparaison puis la mesure.}
\champ{Compétences travaillées}{\textbf{6C04-GEO10} -- Nommer un angle.\newline \textbf{6C04-GEO11} -- Connaître le vocabulaire d'un angle : sommet et côtés.\newline
Compétences mathématiques : \textit{Modéliser} (passer d'une porte, d'un champ de vision à deux demi-droites), \textit{Représenter} (coder un angle par un arc), \textit{Communiquer} (vocabulaire et notations : $\widehat{ABC}$, $[BA)$).\newline
Socle : domaine 1 (langages mathématiques), domaine 2 (méthodes : observer, comparer, ranger).}
\champ{Objectifs de la séance}{À la fin de la séance, l'élève est capable :
\begin{itemize}[nosep,leftmargin=1.2em]
\item de \textbf{reconnaître un angle} dans une situation concrète comme une \textbf{ouverture} (porte, champ de vision) ;
\item de dire qu'un angle est formé par \textbf{deux demi-droites de même origine} ;
\item de donner le \textbf{sommet} et les \textbf{côtés} d'un angle avec la bonne notation ($[BA)$, sommet en premier) ;
\item de \textbf{nommer un angle} avec trois lettres, le sommet au milieu ($\widehat{ABC}$ ou $\widehat{CBA}$).
\end{itemize}}
\champ{Prérequis}{Point, demi-droite et sa notation $[Ox)$ (séquence 2) -- rappelés en tête de l'activité ; repérer un point sur une figure.}
\champ{Matériel}{\textbf{Professeur} : vidéoprojecteur ; diaporama \textbf{Act-6C04-03} (calcul mental, défilement automatique) ; fiche \textbf{Act-6C04-01} projetée ; une vraie porte de la salle pour la démonstration ; craies de couleur.\newline
\textbf{Élèves} (pas d'ardoise) : \textbf{cahier d'exercices} pour tout le travail de recherche (calcul mental, activité, exercices), cahier de cours, colle ; fiches \textbf{Act-6C04-01} (découverte), \textbf{Ex-6C04-01} (nommer un angle) et \textbf{Ex-6C04-02} (sommet et côtés), corrigés imprimés (page 2 de chaque fiche) au bureau.}
\champ{Trace écrite}{Cours, partie I : définition, sommet et côtés (figure légendée), notation $\widehat{ABC}$, exemple et remarque ($\widehat{BAC}$ n'est pas le même angle). Introduite par le « Je retiens » de l'activité.}
\end{tabularx}

\vspace{2mm}
\textbf{Obstacles prévisibles :} croire que la taille de l'angle dépend de la \textbf{longueur des côtés} (une porte plus large n'est pas plus ouverte : ouverture de la porte de la salle à mi-course, puis comparaison avec la porte de l'activité) ;
mettre le sommet en \textbf{première ou dernière} lettre ($\widehat{BAC}$ pour $\widehat{ABC}$) ; oublier le chapeau ;
écrire un côté $[AB)$ au lieu de $[BA)$, ou $[BA]$ (segment) : les côtés ne s'arrêtent pas ;
plusieurs angles dans une même figure (triangle, quadrilatère) : repérer l'angle marqué par l'arc ;
vocabulaire nouveau en grand nombre pour des élèves de REP+ : peu de mots, toujours les mêmes, écrits au tableau.

\newpage

\section*{Déroulement}

\renewcommand{\arraystretch}{1.12}
\small
\begin{tabularx}{\linewidth}{|>{\raggedright\arraybackslash}p{2.1cm}|>{\centering\arraybackslash}p{0.9cm}|X|X|>{\raggedright\arraybackslash}p{2cm}|}
\hline
\textbf{Phase} & \textbf{Durée} & \textbf{Rôle du professeur} & \textbf{Activité des élèves} & \textbf{Modalité / supports} \\ \hline

1. Entrée en classe, automatismes &
\minutes{10} &
Lance le diaporama Act-6C04-03 : 20 questions qui défilent seules, 25 s chacune ($\approx 8$ min) : 10 compléments à $100$ puis 10 à $360$. Correction rapide, fait verbaliser une procédure (« $360 - 270$ : je pense $270 + 30 = 300$, $+ 60$ »). Annonce : « $360$ reviendra dans ce chapitre ». &
Écrivent la date et « Calcul mental » dans le \textbf{cahier d'exercices}, numérotent de 1 à 20 et répondent sans poser l'opération. Corrigent au stylo vert. &
Individuel puis collectif. Diaporama Act-6C04-03, cahier d'exercices. \\ \hline

2. Lancement : la porte de la classe &
\minutes{3} &
Annonce le nouveau chapitre « Les angles ». Ouvre la porte de la salle un peu, puis beaucoup : « Qu'est-ce qui change ? » Fait émerger le mot \textbf{ouverture}. Rappelle la demi-droite (encadré « Rappel » de la fiche). &
Répondent à l'oral, un élève rappelle ce qu'est une demi-droite et sa notation $[Ox)$. &
Collectif, oral. \\ \hline

3. Activité de découverte &
\minutes{12} &
Distribue la fiche Act-6C04-01. Consigne orale courte, une partie à la fois. Partie 1 (portes) : laisse chercher 2 min. Partie 2 (champs de vision) : fait observer la partie grise, « ce que l'animal voit sans bouger les yeux ». Circule : relance les élèves bloqués sur le rangement (d.) en faisant comparer deux animaux à la fois. &
Partie 1 : entourent la porte la plus ouverte et complètent la phrase.\newline Partie 2 : répondent aux questions a, b, c puis rangent les 7 animaux du plus petit au plus grand champ de vision. &
Individuel, entraide en binôme. Fiche Act-6C04-01 collée dans le cahier d'exercices. \\ \hline

4. Mise en commun &
\minutes{5} &
Correction projetée ; un élève vient montrer au tableau le rangement. Fait formuler : la porte et le mur, les deux bords de la zone grise, ce sont \textbf{deux demi-droites qui partent du même point} ; c'est l'\textbf{ouverture} qui compte. Fait compléter « Je retiens ». Trace au tableau deux demi-droites de même origine et nomme : \textbf{angle}, \textbf{sommet}, \textbf{côtés}. &
Expliquent leur rangement, complètent « Je retiens » (angles, demi-droites). &
Collectif, oral. \\ \hline

5. Trace écrite &
\minutes{15} &
Écrit la partie I du cours au tableau ; les élèves copient au fur et à mesure. Définition, figure légendée (sommet, côtés $[BA)$, $[BC)$), notation $\widehat{ABC}$ : « la lettre du milieu est le sommet ». Fait dire à voix haute le nom de l'angle puis ses côtés. Termine par l'exemple et la remarque ($\widehat{BAC}$ faux). Circule pour vérifier les figures et les chapeaux. &
Copient dans le cahier de cours : titre du chapitre, définition, figure légendée, notation, exemple, remarque. &
Collectif. Tableau (craie de couleur pour le sommet), cahier de cours. \\ \hline

6. Exercices &
\minutes{8} &
Présente l'exemple grisé de la fiche Ex-6C04-01 (exercice 1, a.) : « je cherche l'arc, le sommet va au milieu ». Circule : aide prioritaire aux élèves repérés en phase 3. Valide puis donne le \textbf{corrigé} (page 2) pour l'autocorrection. &
\textbf{Tous} : fiche Ex-6C04-01, exercices 1 et 2 (nommer un angle), autocorrection au stylo vert.\newline
\textbf{Élèves rapides} : exercices 3 et 4 de la fiche, puis fiche Ex-6C04-02 (sommet et côtés). &
Individuel. Fiches collées dans le cahier d'exercices, corrigés au bureau. \\ \hline

7. Bilan et devoirs &
\minutes{2} &
Question flash orale : « Dans $\widehat{RST}$, quel est le sommet ? » Donne le travail à la maison. &
Répondent ; notent les devoirs : apprendre la partie I du cours, finir les exercices 1 et 2 d'Ex-6C04-01. &
Collectif, cahier de textes. \\ \hline

\end{tabularx}
\normalsize

\newpage
\renewcommand{\arraystretch}{1.35}
\begin{tabularx}{\linewidth}{|>{\raggedright\arraybackslash}p{3.6cm}|X|}
\hline
\champ{Différenciation}{\textbf{Aide} : rappel de la demi-droite en tête de l'activité, consignes d'une ligne, réponses à compléter sur pointillés ; exemple traité (grisé) en tête de chaque exercice ; rappel de leçon encadré sur les fiches ; cours sous les yeux ; entraide en binôme.\newline
\textbf{Approfondissement} : exercices 3 et 4 de la fiche Ex-6C04-01 (plusieurs angles dans une figure, plusieurs noms pour un angle) ; fiche Ex-6C04-02 (sommet et côtés) ; rôle de tuteur.}
\champ{Évaluation}{\textbf{Diagnostique / formative} : calcul mental (cahier d'exercices), réponses au rangement des animaux (phase 3), question flash en fin de séance, autocorrection des fiches, relevé en passant dans les rangs.\newline
\textbf{Sommative} (séance 3) : DS-6C04-01 (nommer un angle au recto, sommet et côtés au verso), noté sur 40 avec grille de compétences ; QCM bilan DS-6C04-05.}
\champ{Prolongement}{Séance 2 : comparer des angles \textbf{sans les mesurer} (gabarits, Act-6C04-02), puis angles particuliers (nul, aigu, droit, obtus, plat) -- \textbf{Prep-6C04-02}. On y réinvestit l'idée de la séance 1 : la taille d'un angle ne dépend pas de la longueur de ses côtés.}
\champ{Bilan \textit{a posteriori}}{\rule{0pt}{5cm}}
\end{tabularx}


\end{document}

%%% Local Variables:
%%% mode: LaTeX
%%% TeX-master: t
%%% End:
